\documentclass{article}
\usepackage{amsmath,amssymb}
\usepackage[margin=1in]{geometry}
\usepackage{xurl}
\usepackage[hidelinks]{hyperref}
\setlength{\emergencystretch}{2em}
% Shared commands for notes in docs/paper-gaps/.

\DeclareUnicodeCharacter{2080}{\ifmmode{}_{0}\else\textsubscript{0}\fi}
\DeclareUnicodeCharacter{2081}{\ifmmode{}_{1}\else\textsubscript{1}\fi}
\DeclareUnicodeCharacter{2082}{\ifmmode{}_{2}\else\textsubscript{2}\fi}
\DeclareUnicodeCharacter{2099}{\ifmmode{}_{n}\else\textsubscript{n}\fi}

\newcommand{\Lean}{\textsc{Lean}}
\ExplSyntaxOn
\NewDocumentCommand{\leanid}{m}{
  \ifmmode
    \textsf{\detokenize{#1}}
  \else
    \group_begin:
    \urlstyle{sf}
    \tl_set:Nn \l_tmpa_tl {#1}
    \tl_replace_all:Nnn \l_tmpa_tl {\_} {_}
    \exp_args:NV \path \l_tmpa_tl
    \group_end:
  \fi
}
\ExplSyntaxOff
\providecommand{\tr}{\operatorname{tr}}
\newcommand{\tnleanIssueBase}{https://github.com/LionSR/TNLean/issues/}
\newcommand{\tnleanPrBase}{https://github.com/LionSR/TNLean/pull/}
\newcommand{\tnleanDocBase}{https://sirui-lu.com/TNLean/docs/}
\newcommand{\ghissue}[1]{\href{\tnleanIssueBase#1}{GitHub issue \##1}}
\newcommand{\ghpr}[1]{\href{\tnleanPrBase#1}{PR \##1}}
\newcommand{\doclink}[1]{\href{\tnleanDocBase#1.html}{\path{#1}}}

% Dirac notation and the identity operator, as in blueprint/src/macros/common.tex.
% \providecommand keeps note-local definitions authoritative.
\usepackage{dsfont}
\providecommand{\ket}[1]{|#1\rangle}
\providecommand{\bra}[1]{\langle #1|}
\providecommand{\braket}[2]{\langle #1 | #2 \rangle}
\providecommand{\ketbra}[2]{|#1\rangle\!\langle #2|}
\providecommand{\Id}{\mathds{1}}
\providecommand{\one}{\mathds{1}}
\providecommand{\C}{\mathbb{C}}
\providecommand{\MN}[1]{M_{#1}(\C)}
\providecommand{\MD}{M_D(\C)}

% Verdict marker: \gapnote{<kind>}{<status>}. Kind is the policy
% classification (clarification, local-correction, scope-restriction,
% unfaithful, false-source, open-gap); status is open, wip (elimination
% underway), resolved, or historical. Machine-read by the site index;
% prints a verdict line.
\newcommand{\gapnote}[2]{\par\noindent\textbf{Verdict:} #1 (#2).\par}

\title{The Virtual Unitary in the Symmetric Isometric Deformation}
\author{}
\date{2026-10-02}
\begin{document}
\maketitle
\gapnote{scope-restriction}{wip}

\section{The source assertion}
Schuch, P\'erez-Garc\'ia, and Cirac \cite{Schuch2011Phases} show that the
isometric deformation can preserve an imposed on-site symmetry.
The local source is \path{Papers/1010.3732/paper_v3.tex}, lines 645--676.
Their argument first puts the tensor map $P$ in a standard virtual gauge,
then associates to a physical symmetry $U$ a virtual unitary $X$ with
$UP=PX$. Writing $P=QW$, they deduce $[U,Q]=0$ and preservation of the
symmetry by $P_\gamma=Q_\gamma W$ and its parent Hamiltonian.

\section{The proved polar implication}
For any rectangular matrix $P$ and supplied unitaries $U,X$ satisfying
$UP=PX$, one has $[U,PP^*]=0$. Let $E$ project onto $\operatorname{ran}P$.
The positive factor used on the original physical space is
\begin{align}
    Q&=(PP^*)^{1/2}+I-E,
    \label{eq:symmetric_polar_factor}\\
    P&=QW.
    \label{eq:symmetric_polar_identity}
\end{align}
Functional calculus gives $[U,(PP^*)^{1/2}]=[U,E]=0$, hence $[U,Q]=0$.
Since $Q>0$, cancellation in $UP=PX$ gives $UW=WX$.
Thus $Q_\gamma=\gamma Q+(1-\gamma)I$ commutes with $U$, and
$UQ_\gamma W=Q_\gamma WX$.
These statements are proved in
\path{TNLean/MPS/ParentHamiltonian/IsometricDeformationCovariance.lean}.
The identity extension in (\ref{eq:symmetric_polar_factor}) is included:
its square is not identified with $PP^*$ outside $\operatorname{ran}P$.

\section{The parent-space implication}
Suppose in addition that the original two-site MPS boundary space is
invariant under $U\otimes U$. Invertibility of $Q_\gamma$ and commutation
with $U$ preserve this invariance along the deformation. The orthogonal
parent projection then commutes with $U\otimes U$. The source's positive
interaction
\begin{align}
    h_\gamma
        &=(Q_\gamma^{-1})^{\otimes2}h_0(Q_\gamma^{-1})^{\otimes2}
    \label{eq:symmetric_polar_interaction}
\end{align}
has the same commutation property. Each embedded local interaction
commutes with $U^{\otimes N}$, so both open and periodic interaction sums
preserve the symmetry.
These conditional statements are proved in
\path{TNLean/MPS/ParentHamiltonian/IsometricDeformationSymmetry.lean}.

This implication requires invariance of the original MPS boundary space.
An arbitrary virtual unitary $X$ satisfying $UP=PX$ need not preserve the
virtual contractions defining an MPS. Thus the rectangular covariance
identity by itself is not substituted for the original MPS symmetry.

\section{The proved single-block case}
For a single already-injective block $A$, assume either the trace-preserving
canonical normalization $\sum_i A_i^*A_i=I$ or its unital counterpart
$\sum_i A_iA_i^*=I$, and exact invariance of its periodic
matrix product vectors under a unitary on-site representation. Then
the fundamental theorem supplies a virtual similarity, and uniqueness of
the canonical fixed point makes that similarity unitary up to a scalar.
Conjugation by the resulting virtual unitary gives a unitary action on
matrix coordinates. This derives $UP=PX$ and $[U,Q]=0$ without supplying
a virtual unitary as an additional hypothesis. State symmetry also gives
invariance of the original two-site boundary space. Hence both canonical
parent projections and the positive interactions in
(\ref{eq:symmetric_polar_interaction}) preserve the symmetry throughout
the deformation. The derivation is proved in
\path{TNLean/MPS/ParentHamiltonian/InjectiveIsometricDeformationSymmetry.lean},
and its parent-Hamiltonian consequences in the preceding symmetry module.

The normalization convention requires a coordinate qualification.
For $P_{i,(a,b)}=A_i(a,b)$, the source's
$\operatorname{tr}_{\mathrm{left}}(P^*P)=I_{\mathrm{right}}$ at line 653
is $\sum_i A_i^*A_i=I$. The two proved cases are stated separately in their
respective canonical orientations; the trace-preserving case matches this
source convention. Both assume exact vector symmetry, rather than
invariance only up to a phase.

\section{The remaining source step}
For several normal blocks, the derivation must also retain their symmetry
permutation and show that the resulting action on the joint virtual space
is unitary and preserves the MPS contractions. Until that derivation is
proved, the conditional polar and parent-space implications are marked as
such; they are not counted as a formalization of the unrestricted source
symmetry assertion. The separately established uniform-gap results for
the constructed deformation do not depend on this missing symmetry step.

\bibliographystyle{alpha}
\bibliography{references}
\end{document}
